\section{De The Stack a StarCoder: pipeline de extremo a extremo para datos de código}

Esta sección lleva el filtering protocol general a un code corpus real. La particularidad de los datos de código proviene de las diferencias entre lenguajes, la duplicación de repositorios, los generated files, los secrets, las benchmark solutions, los metadata y los eventos de desarrollo estructurados. Cada paso cambia las tareas que el modelo puede aprender.

\subsection{De raw a training data: una gran reducción no equivale a desperdicio}

The Stack v1 pasa de unos 6.4TB de source code a los cerca de 800GB que usa StarCoder; v2 pasa de unos 32.1TB a unos 3.6TB. Esta reduction acumula la selección de lenguajes, las quality rules, el dedup, la PII y las restricciones de formato.

\lecturefigure{slide-44.jpg}{StarCoder/The Stack v1: de 6.4TB raw a 800GB de training data}{Deck oficial de Loubna Ben Allal, página 44}

\lecturefigure{slide-45.jpg}{StarCoder2/The Stack v2: de 32.1TB raw a 3.6TB de training data}{Deck oficial de Loubna Ben Allal, página 45}

\begin{warningbox}{La tasa de filtrado no es un quality metric}
Eliminar el 90\% puede limpiar duplicates, pero también puede borrar rare languages; conservar el 90\% puede mantener la diversity, pero también puede conservar spam. Hay que reportar los tipos eliminados, la source/language retention, la token distribution y el downstream effect.
\end{warningbox}

\subsection{Language selection y per-language thresholds}

v1 selecciona unos 86 de 358 languages y excluye configs y lenguajes sin mantenimiento; v2 se amplía a 600+. Estadísticas como la longitud media de línea, el tamaño de archivo o el comment ratio no pueden compartir un mismo umbral, porque cada lenguaje tiene idioms y generated-code patterns distintos.

\lecturefigure{slide-46.jpg}{The Stack v1: selección de lenguajes, muestreo comunitario y per-language filtering}{Deck oficial de Loubna Ben Allal, página 46}

\begin{knowledgebox}{La forma concreta del humano en el ciclo}
La comunidad BigCode revisa unas 100 muestras por extension, registra el código real, los generated/vendor data, las longitudes y los formatos anómalos, y a partir de ello deriva heuristics. No se trata de que los voluntarios aprueben archivo por archivo, sino de calibrar scalable rules con domain knowledge.
\end{knowledgebox}

\teachervoice{00:30:00--00:33:20, el ponente enfatiza que un mismo average-line-length threshold no puede usarse para todos los lenguajes; la inspection de la comunidad proporcionó el fundamento para diseñar umbrales por lenguaje.}

\subsection{Near-dedup: por qué MinHash y LSH se adaptan a TB-scale code}

El exact dedup solo elimina archivos idénticos; el near-dedup debe identificar cambios de nombre, cambios de formato o modificaciones locales. Si cada archivo se representa como un conjunto de shingles $A,B$, la Jaccard similarity es
\[
J(A,B)=\frac{|A\cap B|}{|A\cup B|}.
\]
MinHash cumple que la probabilidad de colisión del hash es aproximadamente $\Pr[h(A)=h(B)]=J(A,B)$; después, LSH (Locality-Sensitive Hashing) coloca las signatures similares en cubetas candidatas, lo que evita la comparación de todos contra todos.

\lecturefigure{slide-47.jpg}{Near-deduplication: MinHash+LSH y la ganancia en la ablation de StarCoder}{Deck oficial de Loubna Ben Allal, página 47}

\begin{importantbox}{Por qué el near-dedup es un filter de alto valor}
Es independiente del lenguaje, fácil de paralelizar y reduce directamente el entrenamiento repetido y la memorization. La ablation de la slide muestra que pass@1 mejora claramente en varios code subsets, por lo que el equipo adoptó un strong dedup; esto es más confiable que la heuristic de stars/comments sin validar.
\end{importantbox}

\teachervoice{00:33:20--00:35:20, el ponente dice que el near-dedup es el filter con mayor efecto y que no requiere modificarse por separado para cada uno de los 86 lenguajes; cumple a la vez con la ganancia y la escalabilidad.}

\subsection{PII, secrets y benchmark decontamination}

Ni siquiera el secret scanning integrado de GitHub garantiza que el corpus esté libre de emails, names, passwords y keys. BigCode primero anotó manualmente PII spans y luego entrenó el StarPII NER model para escanear la totalidad de los datos; en clase se informó que este paso consumió unas 800 GPU hours. Después hay que retirar también los evaluation benchmarks para evitar un inflated score.

\lecturefigure{slide-48.jpg}{The Stack: PII removal y benchmark decontamination}{Deck oficial de Loubna Ben Allal, página 48}

\begin{knowledgebox}{Primera definición: PII y decontamination}
La PII (Personally Identifiable Information) incluye names, emails, phones y otros datos que permiten identificar a una persona; los secrets incluyen además tokens/passwords/keys. La \textbf{Decontamination} consiste en retirar de los datos de entrenamiento los benchmark items y sus versiones aproximadas o derivadas, para proteger la evaluation independence.
\end{knowledgebox}

\begin{warningbox}{Un Detector solo reduce el riesgo; no puede certificar «cero PII»}
El NER tiene false negatives, las cadenas de código pueden estar ofuscadas y un commit histórico puede filtrar contenido eliminado. Al publicar se necesitan además sample audits, un reporting channel, gating e incident response.
\end{warningbox}

\teachervoice{00:34:00--00:36:20, el ponente describe con mucho detalle la annotation, el NER y el costo de 800 GPU-hour del tratamiento de la PII, y también recuerda que el evaluation set debe retirarse del training corpus.}

\subsection{Formatting: convertir el repository context en señal de entrenamiento}

Después del filtrado aún hay que decidir el sequence format. StarCoder agrega metadata tokens de repository/name/file antes del code; los issues incluyen title, comment y user roles. El formato permite que el modelo aprenda los límites del contexto y también define una interfaz que puede controlarse durante la inferencia.

\lecturefigure{slide-49.jpg}{StarCoder data formatting: metadata de repository, file, stars e issue}{Deck oficial de Loubna Ben Allal, página 49}

\lecturefigure{slide-50.jpg}{Formato estructurado de Git commits y Jupyter notebooks}{Deck oficial de Loubna Ben Allal, página 50}

Para el fill-in-the-middle (FIM), el documento se divide en prefix $P$, middle $M$ y suffix $S$, y la secuencia de entrenamiento puede reordenarse como
\[
\langle\text{PRE}\rangle P\langle\text{SUF}\rangle S
\langle\text{MID}\rangle M,
\]
se sigue usando el next-token loss, pero el modelo aprende a generar el code intermedio cuando conoce el contexto anterior y el posterior.

\begin{warningbox}{Los Metadata tokens también pueden convertirse en un shortcut}
Si las stars, el repository name o la user identity están muy correlacionados con las etiquetas de calidad, el modelo puede memorizar el source bias. Hay que evaluar la capacidad tras retirar los metadata, el privacy risk y la generalization entre repositorios.
\end{warningbox}

\subsection{StarCoder2 mixture: cada source aporta tareas distintas}

La tabla de datos de entrenamiento enumera The Stack, GitHub issues, pull requests, notebooks, documentation, math/code datasets, etc., e indica si se aplica FIM, las epochs y el sample weight. Al leer la tabla no basta con mirar el porcentaje de tokens; hay que preguntarse qué behavior enseña cada tipo de source.

\lecturefigure{slide-51.jpg}{StarCoder2 training-data mixture y source-specific settings}{Deck oficial de Loubna Ben Allal, página 51}

Si la source $i$ tiene $T_i$ tokens y un sampling weight $w_i$, la probabilidad real de muestreo es
\[
p_i=\frac{w_iT_i}{\sum_j w_jT_j}.
\]
Una source pequeña puede someterse a oversample con $w_i>1$, pero la exposure repetida también aumenta la memorization, por lo que debe calibrarse con validation/evaluation.

\begin{knowledgebox}{Source-to-behavior mapping}
Los Code files enseñan syntax/algorithms; los issues enseñan problem description y discussion; los commits/PRs enseñan edit rationale; los notebooks enseñan el entrelazado code--output--markdown; la documentation enseña API usage; los math datasets enseñan symbolic/reasoning. La Mixture es diseño de capacidades, no solo concatenación.
\end{knowledgebox}

\subsection{Tooling: un pipeline de datos reproducible no puede esconderse en un notebook de un solo uso}

La última página enumera datasets, datatrove, datatrove-like distributed pipelines, tokenizers/training code y el BigCode processing repo. El TB-scale filtering requiere streaming, sharding (dividir archivos o muestras en fragmentos que distintos workers pueden procesar de forma independiente), checkpointing, deterministic transforms y dataset versioning.

\lecturefigure{slide-52.jpg}{Tooling de filtrado de datos y entrenamiento: datasets, datatrove, frameworks de entrenamiento y BigCode code}{Deck oficial de Loubna Ben Allal, página 52}

\begin{lstlisting}[caption={Contrato por etapas de un pipeline de datos reproducible}]
discover_sources_and_record_provenance()
normalize_and_language_classify()
apply_quality_filters_with_versioned_thresholds()
exact_and_near_deduplicate()
detect_pii_secrets_and_benchmark_overlap()
format_sequences_and_sample_mixture()
publish_manifest_hashes_stats_and_opt_out_state()
\end{lstlisting}

\subsection{Resumen de la sección}

La data advantage de StarCoder no proviene de un filter mágico, sino de la combinación de source selection, per-language inspection, near-dedup, PII/decontamination, structured formatting, mixture y reproducible tooling. Cada paso tiene costos, modos de falla y evidence.
